\providecommand{\Ell}{\mathrm{Ell}}
\providecommand{\Ch}{\mathrm{ch}}
\providecommand{\BKM}{\mathrm{BKM}}
\providecommand{\EZ}{\mathrm{EZ}}
\providecommand{\KKK}{\mathrm{K3}}

\section*{The two normalisations of $\phi_{0,1}$ and the K3 elliptic genus}

\paragraph{Setting.}
Two normalisations of the weight-$0$ index-$1$ weak Jacobi form coexist
in the literature, and the factor of $2$ between them is the factor of
$2$ between the topological Euler characteristic of K3 and the constant
term of $\phi_{0,1}$. The distinction is structural, not a typo.
Write
\[
  \phi_{0,1}^{\EZ}(\tau, z)
    \;=\; \zeta^{-1} + 10 + \zeta
      \;+\; q\bigl(10\zeta^{-2} - 64\zeta^{-1} + 108 - 64\zeta + 10\zeta^{2}\bigr)
      \;+\; O(q^{2}),
  \qquad \zeta = e^{2\pi i z},
\]
the Eichler--Zagier (1985) normalisation: unique up to scalar in
$J_{0,1}^{\mathrm{weak}}$ with constant term $\zeta^{-1}+10+\zeta$.
Its specialisation at $z=0$ is
\[
  \phi_{0,1}^{\EZ}(\tau, 0)
    \;=\; 1 + 10 + 1 \;=\; 12 \;+\; O(q).
\]
The K3 elliptic genus in the Kawai--Yamada--Yang (1994) convention
satisfies $\Ell(\KKK; \tau, 0) = \chi_{\mathrm{top}}(\KKK) = 24$,
forcing
\[
  \boxed{\;\phi_{0,1}^{\KKK}(\tau, z)
    \;:=\; \Ell(\KKK; \tau, z)
    \;=\; 2\,\phi_{0,1}^{\EZ}(\tau, z).\;}
\]

\paragraph{Attack cycle 1 ($z=0$ specialisation).}
Eichler--Zagier Theorem~9.3 fixes the $q^{0}$-expansion
$\phi_{0,1}^{\EZ}\!\mid_{q^{0}} = \zeta^{-1} + 10 + \zeta$ uniquely.
Setting $\zeta=1$ collapses $\zeta^{-1}+10+\zeta$ to $12$, not $24$.
The K3 $\hat A$-genus integral $\int_{\KKK}\hat A(TX)\wedge\mathrm{ch}(y\cdot TX)\!\mid_{y=1}
= \chi_{\mathrm{top}}(\KKK) = 24$ is the Hirzebruch--Riemann--Roch
output at the Witten-index point. The factor $24/12 = 2$ is then
forced.

\paragraph{Attack cycle 2 (which source says what).}
Slide 919 (Lorgat 2020 presentation deck) writes
$\Ell(\KKK) = 2\phi_{0,1}$, consistent with $\Ell(\KKK; \tau, 0) = 24$
and the Eichler--Zagier normalisation. The April 2020 paper
\texttt{raeez.lorgat.automorphic-corrections.pdf} at \S4, eq.~(4.3)
writes $\Ell(\KKK) = \phi_{0,1}$ -- the factor of $2$ is absorbed
into an implicit redefinition of $\phi_{0,1}$ to the K3 normalisation
$\phi_{0,1}^{\KKK}$, but the paper does not carry the superscript.
The paper is self-consistent under its own convention: its subsequent
line $f(1,1,1)=64$ and $c(0,0)=20$ are K3-normalised values, matching
$c_{\KKK}(0) = 20$ in Eguchi--Ooguri--Tachikawa 2011 Table~1.
The slides are self-consistent under Eichler--Zagier normalisation:
$c_{\EZ}(0,0)=10$, $c_{\EZ}(0,\pm 1)=1$, $\tau(a_{0})=9=10-1$
(Lorgat 2020 \S6 Thm.~4 constant-term balance).
\emph{Both are correct under their own convention; the gap is
normalisation, not content.}

\paragraph{Attack cycle 3 (Chaudhuri--Hockney--Lykken 1995 check).}
Chaudhuri--Hockney--Lykken 1995 \emph{Nucl.~Phys.~B}~459 \S3
computes the K3 elliptic genus from the $\mathcal N=(4,4)$ superconformal
character decomposition of the K3 sigma model and obtains
$\Ell(\KKK; \tau, z) = 2(\zeta^{-1} + 10 + \zeta) + O(q)$, i.e.\
\emph{twice} the Eichler--Zagier form at the leading term. Their
$q^{0}$ reads $2\zeta^{-1} + 20 + 2\zeta$, confirming $\phi_{0,1}^{\KKK} = 2\phi_{0,1}^{\EZ}$
and fixing $\Ell(\KKK; \tau, 0) = 24$ independently of the
Kawai--Yamada--Yang convention.
The same factor appears in Eguchi--Ooguri--Tachikawa 2011 eq.~(2.6),
whose Mathieu-moonshine expansion $\Ell(\KKK) = 24\,\mathrm{ch}_{h=1/4} + \ldots$
has coefficient $24$, not $12$.

\paragraph{Attack cycle 4 ($\kappa_{\BKM} = c_{N}(0)/2$ convention parsing).}
The universal identity $\kappa_{\BKM}(\Phi_{N}) = c_{N}(0)/2$ holds when
$c_{N}(0)$ is the $\zeta^{0}q^{0}$ Fourier coefficient of the
\emph{Eichler--Zagier-normalised} $\phi_{0,1}$ (or its twisted analogue
at $N\ge 2$). At $N=1$:
$c_{1}(0) = c_{\EZ}(0) = 10$, giving $\kappa_{\BKM}(\Delta_{5}) = 5$.
If one substituted the K3-normalised $c_{\KKK}(0) = 20$ naively,
the output would double to $10$, which is $\mathrm{wt}(\Delta_{5}^{2}) = \mathrm{wt}(\Delta_{10})$,
not $\mathrm{wt}(\Delta_{5}) = 5$. The Borcherds weight formula
$\mathrm{wt}(\Phi) = \frac{1}{2}\sum_{\ell}c_{\mathrm{input}}(0,\ell)$
(Borcherds 1998 Thm.~13.3) picks out the \emph{EZ-normalised} input:
the $\frac{1}{2}$ in $c_{N}(0)/2$ and the factor of $2$ in
$\phi_{0,1}^{\KKK} = 2\phi_{0,1}^{\EZ}$ cancel consistently,
so the manuscript convention $c_{1}(0) = 10$, $\kappa_{\BKM} = 5$ is
\emph{invariant under which $\phi_{0,1}$ one calls canonical},
provided one uses the matching normalisation throughout. The
manuscript's choice (EZ, $c_{1}(0) = 10$, $\kappa_{\BKM} = 5$) is the
Eichler--Zagier--Gritsenko--Borcherds lineage.

\paragraph{Attack cycle 5 (downstream propagation check).}
The following load-bearing computations all use the K3-normalised
coefficients $c_{\KKK}(-1) = 2$, $c_{\KKK}(0) = 20$, $c_{\KKK}(3) = 216$,
$c_{\KKK}(4) = -128$, $c_{\KKK}(7) = 1616$, $c_{\KKK}(8) = 1144$:
the $\rho$-stratification PBW recurrence
$d(N) - d(N-1) = |c_{\KKK}(\Delta_{N})|$
(\texttt{chapters/theory/quantum\_groups\_foundations.tex} lines 948--964),
the Lusztig small-form dimensions
$\dim \mathfrak u_{\zeta_{8}}^{\le N} = 8^{2d(N)+3}$
giving $d(1)=2, d(2)=22, d(3)=238, d(4)=366$,
and the Mathieu-moonshine decomposition of $\Ell(\KKK)$.
These all use the factor-of-$2$ correctly: $c_{\KKK}(-1) = 2$
(not $1$) is the K3-normalised polar coefficient, matching
the $2$ real simple roots of $\mathfrak g_{\Delta_{5}}$ and
$\chi_{\mathrm{top}}(\KKK) = 24 = c_{\KKK}(0) + 2 c_{\KKK}(-1) + \ldots$
at the Witten-index specialisation.
\emph{No Wave 11--14 computation is broken by the paper's
implicit-normalisation convention} provided one consistently reads
$\phi_{0,1}$ in the paper as $\phi_{0,1}^{\KKK}$ (K3-normalised).
The slides, writing $\Ell(\KKK) = 2\phi_{0,1}$ explicitly, pick the
EZ normalisation and make the factor of $2$ visible.

\paragraph{Heal (manuscript discipline).}
The Vol~III \texttt{.tex} already carries the superscript discipline
at the point of use: \texttt{quantum\_groups\_foundations.tex}~line~997
states $\phi_{0,1}^{\KKK} = 2\phi_{0,1}^{\EZ}$ with both Eichler--Zagier
and Eguchi--Ooguri--Tachikawa citations, and the downstream BKM
multiplicities use $c_{\KKK}$ superscripted. The Wave-14 Q5 flag is
therefore a drafting-history question about the slide/paper pair,
not a live bug: the Vol~III manuscript is already rectified against
this particular antipattern. The resolution is to adopt the explicit
superscript pair $(\phi_{0,1}^{\EZ}, \phi_{0,1}^{\KKK})$ everywhere a
pre-rectification prose fragment still writes bare $\phi_{0,1}$ in a
factor-of-$2$-sensitive context, and to retain the identity
$\phi_{0,1}^{\KKK} = 2\phi_{0,1}^{\EZ}$ as the bridge.

\paragraph{Summary (three-line verdict).}
Slide $\Ell(\KKK) = 2\phi_{0,1}$: correct in EZ normalisation,
$\phi_{0,1}^{\EZ}\!\mid_{z=0} = 12$, $\Ell(\KKK)\!\mid_{z=0} = 24$.
Paper $\Ell(\KKK) = \phi_{0,1}$: correct in K3 normalisation,
$\phi_{0,1}^{\KKK}\!\mid_{z=0} = 24$, with superscript implicit.
Vol~III manuscript: already uses $\phi_{0,1}^{\EZ}$ and $\phi_{0,1}^{\KKK}$
distinctly; $\kappa_{\BKM}(\Delta_{5}) = c_{1}(0)/2 = 5$ with
$c_{1}(0) := c_{\EZ}(0) = 10$, and downstream uses $c_{\KKK}$
where K3 normalisation is intended. No computation is broken.
